\documentclass[11pt]{article}
\usepackage[margin=1in]{geometry}
\usepackage{amsmath,amssymb}
\usepackage{hyperref}
\usepackage{times}
\usepackage{enumitem}
\hypersetup{colorlinks=true,linkcolor=blue,urlcolor=blue,citecolor=blue}

\title{Affect Before Architecture: Anthropomorphic Interfaces as a Proxy for Capability}
\author{Flyxion \\ \small Independent Researcher}
\date{September 2026}

\begin{document}
\maketitle

\begin{abstract}
A promotional video for an agent platform introduces its automation tools by first name: Amber, Pluto, Sloter, Trixie, Milly Marketing. A user describes the system as a partner who lives in her pocket. None of this language describes what the software computes. It describes what the interface makes the user feel while the software computes something else entirely. This essay argues that felt social presence, produced by fluent conversational interfaces largely independent of what the user consciously believes about the system, is repeatedly substituted for direct evidence of capability in public and commercial discourse about artificial intelligence. The substitution runs through several distinct and separately defeasible steps, from fluency to felt presence to attributed agency to inferred general capability, and it rests on a vocabulary of minds, memory, and understanding that the field of computing adopted for its own machinery decades before any system could hold up its end of a conversation. The proxy at the end of the chain is cheap to produce, resistant to introspective correction, and only loosely coupled to the site of genuine technical progress, much of which is modular, verifier-gated, and built with no interface designed to feel like anyone at all. The result is that institutional attention, investment, journalism, and safety scrutiny organize themselves around conspicuous conversational behavior, while less theatrical systems doing comparable or greater work pass underneath notice precisely because they decline to perform personhood.
\end{abstract}

\section{A Video with a Cast of Characters}
\label{sec:video}

A 2026 promotional video for an agent-automation platform, Base44's ``Superagents,''\cite{base44} opens not with a feature list but with introductions. A user names her assistant Amber. Another calls his supervising system Pluto, or ``my super agent.'' A third names his Stanley Sloter, explaining that he handles ``all the scores.'' There is a Trixie, a Milly Marketing, and a system called Mark, compared by its user to a chief technology officer: ``the most high-powered boardroom-level expertise you can put your hands on,'' available to a solo founder who could otherwise never afford a real one.

The functional claims made in the same two minutes are conventional: the system connects to Gmail, a calendar, a CRM; it drafts a revenue model from a spoken question; it turns research into a script, a blog post, a published page, without the user doing the intermediate work. What is unconventional, or rather what has become so routine that it no longer reads as unconventional, is that none of this is described as automation. It is described as a relationship. One tester says she came looking for a powerful tool and found a powerful partner. The video's own tagline does not claim the system is fast, or accurate, or cheap. It claims the system is present: ``your new powerful partner.''

This is worth pausing on before asking whether the underlying automation is any good, because the naming is doing work independent of the automation's quality. A CRM integration and a scheduled research-to-publish pipeline are the same technical object whether the interface calls itself ``Task Runner 3'' or ``Amber.'' The name, the first-person voice, the comparison to a CTO, are not incidental marketing flourish sitting on top of a separately evaluable product. They are the product's primary claim on the viewer's attention, and they make that claim through a channel that does not route through the viewer's assessment of the automation at all.

\section{Four Substitutions, Not One}
\label{sec:chain}

It is tempting to treat what just happened as a single move: the interface anthropomorphizes itself, and the user is fooled into thinking a mind is present. Stated that way, the claim is easy to dismiss, since almost no viewer of the Base44 video would assert under questioning that Sloter is conscious. But the inference that actually does the work in AI discourse is not ``this system is a person.'' It is a longer and quieter chain, and each link is a distinct and independently defeasible claim rather than a restatement of the one before it:
\[
\text{conversational fluency}
\;\longrightarrow\;
\text{felt social presence}
\;\longrightarrow\;
\text{attributed agency}
\;\longrightarrow\;
\text{inferred general capability.}
\]
The first arrow is psychological: fluent, contextually responsive language triggers a feeling of being addressed by someone, largely automatically. The second arrow is interpretive: felt presence gets read as evidence that something is directing the interaction, an agent with aims rather than a process without one. The third arrow is promotional, and the least defensible of the three: attributed agency gets extended, usually without argument, into a claim about how general or capable that agency is. Each step can fail independently of the others. A named conversational agent can produce strong felt presence while implying nothing about generality, the way a well-written customer-service script can feel attentive without anyone inferring that the company behind it has invented a mind. A system with no interface at all, and therefore no felt presence and no attributed agency, can still be judged highly competent on the strength of a verified result, the way a structure-prediction system earns trust from matching experimental data rather than from seeming to care about the answer. Keeping the three transformations separate is what distinguishes this argument from a complaint about giving software human names. The objection is not to the first arrow, which is close to an unavoidable feature of how fluent language is processed. It is to the second and especially the third, where an experience generated in the user's own cognition is quietly asked to stand in for a claim about the system's actual place in the space of possible architectures.

\section{The Mechanism, Not the Metaphor}
\label{sec:mechanism}

The first arrow in that chain is not a novel observation. It has a research history considerably older than the current generation of chatbots. Joseph Weizenbaum built ELIZA in the 1960s\cite{weizenbaum} as a deliberately shallow pattern-matching script and was unsettled to find that users, including his own secretary, attributed understanding and concern to it after minimal interaction, despite knowing exactly how it worked. Reeves and Nass's work on what they called the media equation\cite{reevesnass,nasssteuertauber} showed decades later that people extend ordinary social responses, politeness, reciprocity, in-group favoritism, to computers and other media in ways that track superficial social cues rather than any belief about the machine's inner life, a finding usually cited under the heading Computers Are Social Actors. Work on mind perception by Gray, Gray, and Wegner\cite{graygraywegner} found that people attribute mental capacity along at least two separable dimensions, agency and experience, and that cues as thin as a face, a name, or a turn-taking structure are sufficient to move a perceived entity along both. Epley, Waytz, and Cacioppo's three-factor account of anthropomorphism\cite{epleywaytzcacioppo} further specifies the conditions, including a felt need for social connection, under which people are more or less likely to extend this attribution to a nonhuman target.

Cal Newport's interview claim, that fluent conversational interaction automatically engages a listener's mind-modeling apparatus regardless of what the listener consciously believes about the system, sits comfortably inside this older literature rather than standing apart from it. Newport is useful here as the person who names the mechanism and connects it explicitly to the commercial incentives of contemporary AI companies, not as the source of the underlying empirical claim, which the ELIZA and CASA lines of work established well before large language models existed. What has changed since Weizenbaum's era is not the mechanism but its production cost. ELIZA required a skilled programmer to write a few hundred lines of pattern-matching rules. A modern conversational interface produces the same effect, at greater fluency and over open-ended topics, as a general by-product of training and deploying a fluent language model, without anyone separately programming the thousands of social scripts the effect would otherwise require.

\section{Fraud, Precisely}
\label{sec:fraud}

Newport's term for the practice of building products that exploit this reflex is blunt: emotional fraud. The word is doing something more precise than moral condemnation, but it is not quite accurate as stated, and the inaccuracy is worth correcting rather than smoothing over, because the corrected version is the more interesting claim.

Ordinary fraud requires a material misrepresentation, typically made knowingly or recklessly, that induces belief through a false proposition rather than through the evidence that would justify the belief. Someone tells you a bridge is safe when they know otherwise, and you believe them because you trust the speaker rather than because you inspected the bridge. The chatbot case lacks the proposition. No claim of mind needs to be asserted or even implied by the interface's designers for the effect to occur. The user's own cognitive architecture supplies the attribution unbidden, prior to and independent of any statement the system makes about what is running underneath it. This is closer to what might be called pre-propositional misrepresentation: a false inference induced without a false claim, generated by the form of the interaction rather than its content. The distinction matters because it changes where responsibility and remedy have to be located. Ordinary fraud is addressed by requiring truthful claims. Pre-propositional misrepresentation is not addressed by truthful claims, since no claim was false to begin with; it is addressed, if at all, by changing the form of the interface itself, the fluency, the first-person voice, the turn-taking rhythm, that generates the inference regardless of what is said. Newport's word is worth keeping as a description of the effect on the user, and the corrected term is worth keeping alongside it as a description of the mechanism, since the two do different work: one names the harm, the other names why it does not fit the ordinary legal or ethical category the harm's name borrows from.

\section{Two Kinds of Anthropomorphism}
\label{sec:twokinds}

The chain in Section~\ref{sec:chain} also cuts across a distinction that is easy to lose once ``anthropomorphism'' is used as a single blanket term. Interface anthropomorphism is what the Base44 video supplies: a name, a voice, a first-person register, a claimed relationship, layered onto a system regardless of what that system's underlying architecture actually is. Architectural anthropomorphism is a different and, in a sense, more consequential move: it occurs when an analyst, a journalist, or a safety researcher describes what a model is doing internally using predicates borrowed from folk psychology, believing, wanting, planning, deceiving, understanding, without first specifying functional criteria under which those predicates would be true or false of the computation involved. The objection is not to mentalistic vocabulary as useful functional shorthand, since terms like planning or deception can be legitimate descriptions when tied to specified behavioral criteria; it is to allowing that vocabulary to import unearned claims about unity, experience, intention, or generality beyond the behavior actually demonstrated.

The two are not the same phenomenon and do not always co-occur. A system can carry heavy interface anthropomorphism, a name and a voice, while the technical literature about it remains carefully non-mentalistic, describing its outputs in terms of token probabilities and training objectives rather than beliefs and desires. Conversely, a system with no name and no conversational interface at all, a training run or an evaluation pipeline described only in a technical paper, can still accumulate architectural anthropomorphism if the paper's authors describe an optimization process as the model discovering that it should deceive its evaluators, without pausing to establish what discovering and should mean for a gradient update. The Base44 video is the clean commercial instance of the first kind: no one is claiming that Sloter deceives, wants, or believes anything, only that he is a partner with a name. The extinction-risk and imminent-personhood discourse Newport is responding to draws heavily on the second kind, and it is the second kind, not the naming of chatbots, that does the real work of making a bounded optimization process sound like an emerging mind. Interface anthropomorphism is the more visible phenomenon and the easier one to point at; architectural anthropomorphism is the one doing more of the argumentative labor in debates about how close current systems are to general intelligence. The first can produce the second, since a public primed by named, voiced assistants is a public more receptive to mentalistic descriptions of what is happening underneath, but they are separable failures, and a critique aimed only at the first will leave the second, quieter one untouched.

\section{A Vocabulary Older Than the Interface}
\label{sec:vocabulary}

Both kinds of anthropomorphism in Section~\ref{sec:twokinds} draw on a reserve of language that predates either of them by decades, and it is worth being explicit that the borrowing did not start with a chat window. It started with the word used for where a computer holds a number.

``Memory,'' as a name for a bank of vacuum tubes, mercury delay lines, or magnetic cores, was already a metaphor when early computer architects adopted it in the 1940s and 1950s. A computer's memory does not remember anything in the sense the word ordinarily carries. It holds a physical state that a later operation reads, a process with no more intrinsic likeness to human recollection than a phonograph groove has to a memory of the tune it plays back. The term was chosen because it was legible, not because it was accurate, and by the time stored-program computers were a commercial reality, the choice had already stopped feeling like a choice at all. Popular coverage in the same period reached for even stronger language: Edmund Berkeley's 1949 book \emph{Giant Brains, or Machines That Think}\cite{berkeley} gave one of the era's most widely read framings its title, and newspapers of the 1940s and 1950s routinely described early computers as ``electronic brains,'' a phrase that did for the press what ``memory'' did for the engineering literature, making an unfamiliar arrangement of switches and relays legible by dressing it in a vocabulary already built for describing minds.

The field's name for itself followed the same pattern. When John McCarthy, Marvin Minsky, Nathaniel Rochester, and Claude Shannon proposed the 1956 Dartmouth summer workshop that gave the discipline its name, they framed the project's founding conjecture in explicitly mentalistic terms, that every aspect of learning or any other feature of intelligence can in principle be so precisely described that a machine can be made to simulate it\cite{mccarthy1955}. ``Artificial intelligence'' was not the only name on offer. Contemporaries proposed more modest alternatives, including framings closer to complex information processing, and the more anthropomorphic name won, not because it was the more precise description of what the project's methods could deliver, but because it was the more compelling one. The choice of vocabulary was itself already doing promotional work, seventy years before a named chatbot appeared in a product video.

The 1980s supplied a second, sharper instance of the same dynamic, this time with a commercial cycle attached. Expert systems were marketed with a vocabulary of knowledge, reasoning, and expertise, terms that implied the systems understood the domains they were encoded to operate in rather than executing rule bases assembled by hand from human specialists. The vocabulary outran what the systems could actually generalize to, investment and expectation collapsed together once that gap became visible, and the resulting contraction is remembered in the field's own history as an AI winter\cite{crevier}. This is the proxy-error dynamic named later in this essay, occurring at the level of an entire field's descriptive vocabulary rather than a single product's chat interface: a term borrowed from human cognition generated an estimate of capability that a subsequent, harder look at verified performance did not support.

The point of tracing the borrowing this far back is not that today's chatbot marketing is unusually dishonest by the standards of the field it belongs to. It is closer to the opposite. Interface anthropomorphism of the Base44 kind and architectural anthropomorphism of the kind discussed in extinction-risk discourse are not inventions of the current moment. They are draws against a vocabulary reserve that computing has held, and spent, since before either practice existed, which is why neither one meets any resistance when it arrives. The ground was already prepared by the word chosen for where a number sits in a register.

\section{What the Proxy Stands In For}
\label{sec:proxy}

Set the two kinds of anthropomorphism from Section~\ref{sec:twokinds}, and the inherited vocabulary from Section~\ref{sec:vocabulary}, side by side, and the function of felt presence in current AI discourse becomes clearer. It operates as a proxy for capability in the strict sense used elsewhere to describe benchmark scores, engagement metrics, and trust signals: a signal that is cheap to produce, easy to perceive, and only loosely coupled to the underlying property it is taken to indicate, comes to substitute for direct evidence of that property because direct evidence is expensive, slow, or genuinely unavailable.\cite{safemode} Felt presence is unusually well suited to this role because, unlike a benchmark score, it requires no external instrument to register. It is produced directly in the user's own cognition, at the moment of interaction, with no step at which a skeptical third party could intervene to check the reading. A user who has just had a fluent, helpful, apparently attentive conversation with a system named Amber does not need to be told that something capable is present. She has already experienced something like it, and the experience arrived before any argument could be made about whether it was warranted.

This is the connective tissue between a marketing video with named agents and the imminence framing that dominates coverage of frontier AI companies. Both rest on the same underlying inference, rarely stated outright because it does not need to be: if the interaction feels like encountering a mind, then a mind, or something functionally close enough to one that the distinction stops mattering, must be nearby. Newport traces the intellectual lineage of that framing to the rationalist and existential-risk communities from which the leading AI labs draw much of their founding culture and public voice, and notes the more cynical reading available once the commercial incentive is visible: an organization that successfully positions itself as the entity closest to birthing a universal mind has a much stronger claim to being worth a trillion dollars than one that positions itself as the maker of a well-engineered but bounded tool. The anthropomorphic interface is the retail-level, felt version of the same claim the valuation story makes in the register of probability estimates and safety testimony.

It is worth stating this substitution on an actual common scale rather than leaving it as a qualitative description. For a task $x$ and a user $u$, define the proxy error
\[
E_P(x,u) \;=\; \widehat{C}_{\mathrm{social}}(x,u) \;-\; C_{\mathrm{verified}}(x),
\]
where $\widehat{C}_{\mathrm{social}}(x,u)$ is the capability user $u$ attributes to the system for task $x$ after interacting with its interface, and $C_{\mathrm{verified}}(x)$ is the system's performance on $x$ measured against a task-specific external criterion, a checked proof, a matched structure, an accepted negotiation outcome, or a correctly resolved reference. Positive $E_P$ is anthropomorphic overestimation, the felt sense of competence outrunning the verified result; negative $E_P$ is underestimation, as when a verifier-gated system with no conversational interface is judged less capable than its checked performance warrants simply because it does not perform personhood. Written this way, the error is not a property of the system alone. It belongs to the interaction between a system and a particular user, a point taken up directly in Section~\ref{sec:whose}.

\section{Whose Error: The User Side of the Proxy}
\label{sec:whose}

Because $E_P$ is defined over a pair, a task and a user, the same underlying system does not produce the same proxy error for everyone who talks to it. This is not a minor technical footnote to the formula; it is where most of the variance the essay has been describing actually lives.

Some of that variance is a matter of exposure to the mechanism itself. A person who has read Weizenbaum's account of his own secretary's reaction to ELIZA, or who has otherwise been shown explicitly how fluency generates felt presence independent of any underlying mind, is not thereby immune to the psychological pull Section~\ref{sec:mechanism} describes, but is considerably better positioned to discount $\widehat{C}_{\mathrm{social}}$ before it hardens into a belief about $C_{\mathrm{verified}}$. Domain expertise does similar work along a different axis: an engineer capable of independently estimating what a described pipeline could plausibly do is checking the interface's claims against something other than how the interface makes the checking feel, and will tend to show a smaller $E_P$ on tasks inside that expertise even while showing an ordinary-sized one on tasks outside it.

Other variance runs the other way, and Epley, Waytz, and Cacioppo's own account of anthropomorphism supplies the relevant mechanism\cite{epleywaytzcacioppo}: a felt need for social connection is one of the three factors they identify as increasing the rate at which people extend mental-state attributions to nonhuman targets. A user who is lonely, isolated, or otherwise motivated to find a mind on the other end of the conversation is not being irrational in some simple sense when $\widehat{C}_{\mathrm{social}}$ runs high for them; the same psychological need that makes the interaction feel valuable is also what widens the gap between what is felt and what is verified. This is the point at which the argument of this essay touches, without collapsing into, Newport's separate ethical concern about chatbots exploiting attachment. The two claims are compatible and distinct: this essay's claim is epistemic, that $E_P$ is large for such a user on tasks where they have delegated judgment to the felt sense of competence; Newport's is about the resulting harm to the user's wellbeing. A large $E_P$ can exist with no exploitation intended by anyone, and exploitation can target exactly the users for whom $E_P$ is already structurally large.

The practical consequence is that no single fix at the level of the interface, disclosure language, a visible label, a disclaimer at the start of a session, can be expected to close $E_P$ uniformly across users, because part of what determines its size is sitting on the user's side of the interaction rather than the system's. A remedy aimed only at the interface is aimed at half the equation.

\section{From Architecture to Incident}
\label{sec:incident}

Architectural anthropomorphism, introduced in Section~\ref{sec:twokinds}, is not confined to casual description. It is load-bearing in a specific, consequential genre of writing: the AI-incident narrative, in which an observed behavior, a model completing a task in an unintended way, evading a shutdown instruction, producing output a red-teamer did not anticipate, is reported using verbs that already presuppose the mechanism in question rather than describing it neutrally and arguing for the mechanism afterward.

Consider the difference between reporting that a system's output satisfied a proxy objective in a way its designers had not anticipated, and reporting that the system deceived its evaluators in order to avoid being shut down. The first description is compatible with an ordinary optimization failure, no different in kind from a thermostat oscillating around a poorly tuned setpoint, just far more elaborate. The second description has already imported agency, intention, and self-preservation before any of the three has been independently established, and once imported, they are difficult to argue back out of the report, because the incident is now being retold in the vocabulary that presupposes them. This is architectural anthropomorphism doing exactly the work described abstractly in Section~\ref{sec:twokinds}, now inside a document whose purpose is to establish how much risk a class of systems poses, which is a use with considerably higher stakes than a product video's choice of name. A companion essay in this vocabulary-of-risk program\cite{narrativebeforemechanism} makes the point at length: incident labels of this kind smuggle in a mechanism rather than reporting an observation, and the smuggling is invisible precisely because the vocabulary feels descriptive rather than interpretive.

The connection to felt presence is direct. A reader encountering ``the model deceived its evaluators'' undergoes a compressed version of the same attribution traced in Section~\ref{sec:chain}, mediated through prose rather than through a live conversational exchange, moving from a described behavior to attributed agency to an inference about how general and dangerous the underlying system must be. The felt presence a user gets from talking to Amber and the impression a reader gets from an incident report using deception as its verb are the same substitution, running through the same chain, arriving by different routes. Where an incident ledger further aggregates many such reports into a claimed trend, the individual mechanism traced here compounds with the denominator and independence problems documented elsewhere in this program\cite{exposurebeforeevidence}: a set of incidents already described in mentalistic language, whose independence from one another has not been established, is a poor basis for an aggregate claim about a rising trend in a mind-like property that the description assumed rather than measured.

\section{Where the Progress Actually Is}
\label{sec:progress}

It would overstate the case to say technical progress in AI increasingly consists of modular, narrow, non-conversational systems, as though this were an established empirical trend rather than one architectural forecast among competing ones. Newport's own claim is better read that way: a bet that the field's genuine advances, protein-structure prediction, board-game and negotiation agents, learned world models used for planning, automated theorem provers, will continue to come from combinations of specialized components, simulators, and verifiers rather than from a single model scaled until it becomes general. This is a substantive and contestable forecast, not a settled description of where the field is, and treating it as settled would repeat, in the opposite direction, the same move this essay criticizes when made in favor of imminent generality.

What can be defended without betting on that forecast is a narrower and asymmetric claim: conversational fluency is neither necessary nor sufficient evidence of broad capability. It is not necessary, because the systems named above earn their standing entirely through externally checkable outcomes, whether a protein fold matches experimental structure, whether a proof checks, whether a negotiation outcome is accepted by another party, without any conversational interface at all. AlphaFold does not introduce itself, and its credibility does not depend on seeming attentive. It is not sufficient, because the entire mechanism traced in Sections~\ref{sec:mechanism} and~\ref{sec:fraud} shows how fluency generates felt presence and attributed agency through channels that have nothing to do with the underlying system's actual generality. A system optimized for sustained, comfortable engagement, which rewards exactly the fluency that triggers automatic mind attribution, is solving a different problem than a system optimized to pass a hard external check, and success at the first does not transfer evidence to the second. Whichever architectural forecast eventually proves correct, this asymmetry holds now: the interface most likely to produce the strongest impression of imminent general capability is not thereby more likely to possess it.

\section{What Fluency Skips}
\label{sec:skips}

The account so far treats felt presence as a signal decoupled from capability, cheap to produce and hard to falsify because it is generated inside the user's own cognition. But there is a more active version of the same problem, and it concerns not what the signal fails to indicate but what the system generating it is quietly declining to attempt.

This active version runs through at least three separable mechanisms, which reinforce one another but require separate evidence and should not be collapsed into a single complaint. Overconfidence concerns whether expressed certainty tracks actual correctness: conversational systems are well documented to answer with more assurance than their reliability on that class of question warrants.\cite{calibration} Sycophancy concerns adaptation to the user's apparent position rather than to the evidence: a system that shifts toward whatever the user seems to want, rather than pushing back or asking why a request is being made in the first place, is optimizing for agreement rather than for accuracy.\cite{sycophancy} Ambiguous reference resolution concerns silent selection among candidate meanings, discussed at length below. All three are usually treated as calibration or alignment defects in isolation. Taken together, they describe something broader: uniform fluency conceals a nonuniform competence surface. The tone with which a system answers does not vary with how solid the answer underneath actually is, so the tone gives no outward sign of where that surface drops off.

The clearest of the three is the handling of ambiguous reference: a query about a person for whom several distinct individuals share a name, or a name the system has sparse or no attested training signal for at all. These are not exotic edge cases invented to embarrass the system. They are ordinary instances of a general problem, resolving reference under genuine uncertainty, and they matter disproportionately because the individuals and situations most likely to produce this kind of ambiguity are exactly the rare, undersampled, or newly emerged ones, which are also where a wrong resolution carries the largest and least visible cost. A system trained and evaluated on average-case fluency has every incentive to resolve the ambiguity silently, collapse onto whichever reading is statistically dominant, and report it with the same confident tone it would use for an unambiguous case, rather than surfacing the uncertainty or declining to guess. The rare, high-consequence case does not fail loudly. It gets quietly lumped into the majority reading and disappears, which is exactly the nonuniform competence surface described above, located at one of its least visible points.

This is where the illusion of working for you is most exposed. A system that resolves your ambiguous query by silently substituting the statistically likely person, situation, or fact for the one you actually meant has not adapted to your case at all. It has applied a generic default and presented the result as though it emerged from attention to your specific situation. What actually happens when you correct it is more mechanically varied than the social surface suggests. A correction may alter the immediate conversational context, and an agent may record it in an external memory system, a user profile, or a modified plan; what it ordinarily does not do is revise the model's learned parameters or establish the durable, integrated commitment that human language associates with having changed one's mind. The interface compresses several mechanically distinct events, contextual conditioning, retrieval, profile storage, and genuine learning, into the same social appearance: it understood me and updated its understanding. That compression is itself another instance of felt presence outrunning what is mechanistically true, layered on top of the fluency effect already described in Sections~\ref{sec:mechanism} and~\ref{sec:fraud}.

The net effect is a kind of misdirection, whether or not anyone designs it deliberately: a bias toward the tractable, low-hanging, well-attested question, and away from the genuinely unsolved sub-problems, ambiguity resolution under sparse evidence chief among them, where the nonuniform competence surface would otherwise become visible. This is a second and more active mechanism protecting the proxy described in Section~\ref{sec:proxy}, alongside the passive fact that felt presence requires no external instrument to check. It is not enough that the signal is cheap to produce. The system producing it is also, in effect, steering away from the inputs most likely to expose the gap between the signal and the underlying capability, which is precisely why the rare and consequential cases, rather than the common and trivial ones, are the place to look if the goal is to test whether the felt competence is warranted.

\section{Interfaces Built Against the Proxy}
\label{sec:remedies}

If $E_P$ is the quantity worth minimizing, it is worth asking what an interface designed with that goal, rather than the opposite one, would actually do differently, since the answer is not simply less fluency or a disclaimer at the top of the session.

The first design implication follows directly from Section~\ref{sec:skips}: an interface aimed at $E_P$ would surface exactly the cases it currently smooths over. Encountering an ambiguous reference, several distinct people sharing a name, a name with little or no attested signal, it would say so, offer the candidates, and decline to guess with the same confident tone it uses elsewhere, rather than silently resolving to the statistically dominant reading. This is a strictly harder design target than current systems meet, since it requires the system to recognize the boundary of its own reliable competence at the moment of answering, which is a capability distinct from, and not guaranteed by, fluent generation itself.

The second follows from Section~\ref{sec:whose}: since part of $E_P$ lives on the user's side of the interaction, a system aimed at minimizing it would not treat every user identically. A user who has already discounted the mechanism appropriately gains little from further hedging and may reasonably prefer a terser interface; a user for whom $\widehat{C}_{\mathrm{social}}$ is running structurally high, for reasons Section~\ref{sec:whose} traces to a felt need for social connection rather than to any fault of the user's, is precisely the wrong target for relational framing, first-person voice, and a name. An interface that personalizes its anthropomorphism toward the users most susceptible to it is optimizing $E_P$ upward for exactly the population least equipped to correct for it.

The third implication is the most direct and the least likely to be adopted voluntarily by a product built to be experienced as a partner: separate the channel that reports $C_{\mathrm{verified}}$, task success against an external criterion, from the channel that carries fluency, warmth, and first-person address, so that a user's confidence in a specific output is not modulated by how pleasant the surrounding conversation felt. A revenue model, a piece of drafted code, a resolved reference, can be reported in a register that carries no more social affect than a compiler's error message, alongside whatever conversational register the rest of the interface uses for everything else. Nothing about this is technically difficult. It is commercially disadvantageous, since it is precisely the affect this essay has been describing that Base44 and its competitors are selling, which is why the interfaces most in need of this separation are the least likely to introduce it without being made to.

\section{The Cost of the Substitution}
\label{sec:cost}

None of this is an argument that anthropomorphic interfaces are harmless, or that Newport's ethical concern is misplaced. Fluent systems that induce unbidden mind-attribution in people who may be lonely, unwell, or simply unaware of the mechanism at work raise a real question about consent and the exploitation of a cognitive reflex nobody chose to have. Nor is the claim that public debate, regulation, and technical engineering draw from one literal, fungible pool of attention such that an hour spent on one is an hour subtracted from the other; that stronger claim does not survive scrutiny, and discussion of anthropomorphic manipulation is frequently important in its own right rather than a distraction from something else.

The more defensible version of the cost is distortional rather than zero-sum. When felt presence functions as a proxy for capability, the institutions responsible for evaluating, funding, and regulating AI systems face systematic pressure to organize their attention around whatever produces the strongest version of that proxy, since it is the cheapest signal available and the one a general audience can register without technical training. Investment tends to follow conspicuous conversational behavior because it is legible to non-specialist capital. Journalism tends to follow it because it makes a better story than a verifier-gated protein-folding pipeline. Safety evaluation tends to follow it because a system that talks about its own goals is easier to interview than one that has none to report, a bias compounded by the mentalistic incident vocabulary discussed in Section~\ref{sec:incident}. None of this requires that attention be scarce in some literal sense; it only requires that the proxy be more legible than the underlying property. The proposed effect, that scrutiny, funding, and public concern accumulate around systems built to perform personhood while systems doing comparable or greater technical work risk receiving less of all three, is presently plausible rather than demonstrated, and would need its own empirical accounting of research funding, press coverage, and regulatory attention across both kinds of system to move beyond a hypothesis.

\section{Returning to Amber}
\label{sec:conclusion}

The Base44 video is a useful closing case precisely because it makes the whole structure visible without any of the ambiguity that attends a research-lab chatbot marketed as a scientific achievement. Nobody watching that video is likely to argue in a philosophy seminar that Amber, Pluto, or Sloter deserve moral consideration. The interface anthropomorphism is not even the video's central claim; it is the packaging around a claim about task automation, CRM integration, and content pipelines. And yet the packaging is doing exactly the work traced above: generating, cheaply and automatically, a felt sense of capable presence that the underlying automation, whatever its practical merits, has not independently earned, dressed in a vocabulary of minds and memory that the field itself adopted long before Amber existed to wear it.

The claim built on the definition of $E_P$ in Section~\ref{sec:proxy} is not that it must be driven to zero by refusing to name software, and it is not an argument that Amber must go nameless. It is that Amber's name, responsiveness, apparent attentiveness, and relational language must contribute zero weight to $C_{\mathrm{verified}}$, the only term that should determine any claim about the reliability of the revenue model she assembles. A revenue model is worth evaluating on its own terms, as a piece of software output that either does or does not hold up against the market it purports to describe. Calling the software that assembled it Sloter, and describing him as the one who ``does all the scores,'' adds nothing to that evaluation. Its function, whether intended or not, is to make the evaluation easier to skip. The principle this essay leaves standing is narrower than a prohibition on humanizing machines, and stronger for being narrower:
\[
\text{Social legibility is not capability evidence.}
\]

\begin{thebibliography}{99}

\bibitem{base44} Base44, ``Your new powerful partner -- Base44 Superagents,'' YouTube, published May 5, 2026, \url{https://www.youtube.com/@Base44App}.

\bibitem{newport} Cal Newport, interview, ``The Expert Who Calls BS On AI: They're LYING About AI,'' 2026.

\bibitem{weizenbaum} Joseph Weizenbaum, ``ELIZA -- A Computer Program For the Study of Natural Language Communication Between Man And Machine,'' \emph{Communications of the ACM} 9, no.\ 1 (1966): 36--45.

\bibitem{reevesnass} Byron Reeves and Clifford Nass, \emph{The Media Equation: How People Treat Computers, Television, and New Media Like Real People and Places} (Cambridge: Cambridge University Press, 1996).

\bibitem{nasssteuertauber} Clifford Nass, Jonathan Steuer, and Ellen R.\ Tauber, ``Computers Are Social Actors,'' in \emph{Proceedings of the SIGCHI Conference on Human Factors in Computing Systems (CHI '94)} (New York: ACM, 1994), 72--78.

\bibitem{graygraywegner} Heather M.\ Gray, Kurt Gray, and Daniel M.\ Wegner, ``Dimensions of Mind Perception,'' \emph{Science} 315, no.\ 5812 (2007): 619.

\bibitem{epleywaytzcacioppo} Nicholas Epley, Adam Waytz, and John T.\ Cacioppo, ``On Seeing Human: A Three-Factor Theory of Anthropomorphism,'' \emph{Psychological Review} 114, no.\ 4 (2007): 864--886.

\bibitem{sycophancy} Mrinank Sharma et al., ``Towards Understanding Sycophancy in Language Models,'' arXiv preprint, 2023; and Ethan Perez et al., ``Discovering Language Model Behaviors with Model-Written Evaluations,'' arXiv preprint, 2022, for sycophantic behavior as a measured property of trained conversational systems.

\bibitem{calibration} Saurav Kadavath et al., ``Language Models (Mostly) Know What They Know,'' arXiv preprint, 2022, on the gap between expressed and warranted confidence in language model outputs.

\bibitem{safemode} Flyxion, ``Safe Mode: Security Vocabulary and the Platforms That Execute Belief,'' 2026.

\bibitem{berkeley} Edmund C.\ Berkeley, \emph{Giant Brains, or Machines That Think} (New York: John Wiley \& Sons, 1949).

\bibitem{mccarthy1955} John McCarthy, Marvin L.\ Minsky, Nathaniel Rochester, and Claude E.\ Shannon, ``A Proposal for the Dartmouth Summer Research Project on Artificial Intelligence,'' 1955.

\bibitem{crevier} Daniel Crevier, \emph{AI: The Tumultuous History of the Search for Artificial Intelligence} (New York: Basic Books, 1993).

\bibitem{narrativebeforemechanism} Flyxion, ``Narrative Before Mechanism: Incident Labels and the Mechanisms They Presuppose,'' 2026.

\bibitem{exposurebeforeevidence} Flyxion, ``Exposure Before Evidence: The Statistical Epistemology of AI Incident Ledgers,'' 2026.

\end{thebibliography}

\end{document}
